\section{The simulation core}\label{sec:sim}

The simulation core is where the codebase's central habit --- build the seam first, wire in the
consumer later --- is most visible. What actually flies today is a three-degree-of-freedom point
mass: the integrated state is center-of-mass position and velocity, nothing more. Around that
modest physics sits a genuinely pluggable machine: two interchangeable ODE solver backends selected
by name at runtime, a string-keyed registry of three atmosphere and two gravity models, a
hang-proof run loop that maps every possible exit onto a typed verdict, and per-step trajectory
statistics. The scaffolding for more --- quaternion attitude state, a torque interface, per-step
inertia recording, a wind model class --- is present, typed, and in several cases tested, but has
no consumer. This section reviews what is integrated, how the loop terminates, how the two solvers
differ, what the environment registry offers, and how the Qt-free \code{QtRocket} controller
drives it all.

\subsection{The injected ODE and what is actually integrated}\label{sec:sim-ode}

The physics enters the solver as a closure built once in the \code{Propagator} constructor
(\src{sim/Propagator.cpp}{36}):

\begin{minted}{cpp}
std::function<std::pair<Vector3, Vector3>(double, Vector3&, Vector3&)> linearODEs = [this](double t, Vector3& state, Vector3& rate) -> std::pair<Vector3, Vector3>
{
    Vector3 dPosition;
    Vector3 dVelocity;
    dPosition = rate;
    dVelocity = object->getForces(t, state, rate, *environment) / object->getMass(t);

    return std::make_pair(dPosition, dVelocity);
};
\end{minted}

\noindent This lambda is the entire integrated physics: $\dot{x}=v$, $\dot{v}=F(t,x,v)/m(t)$ over
a coupled pair of \code{Vector3}s. Time is threaded through deliberately, so that the motor thrust
curve is sampled at each Runge--Kutta stage's node time rather than frozen at the step start
(\src{sim/Propagator.cpp}{34}, \src{sim/DESolver.h}{41}). The \code{model::Propagatable} behind
\code{object} is in practice \code{RocketModel}, the model$\leftrightarrow$sim bridge described in
\cref{sec:model}.

\begin{designrule}[Environment by injection, not by singleton]
The \code{Propagator} holds the \code{Environment} and passes \cppinline|*environment| into every
\code{getForces} call (\src{sim/Propagator.h}{106}, \src{sim/Propagator.cpp}{41}). The force path
never reaches back to a global for its physics models, which is what lets tests swap atmospheres
and gravity per-run and keeps the \code{utils}$\rightarrow$controller back-call cycle out of the
dependency graph.
\end{designrule}

The force model itself (\src{model/RocketModel.cpp}{68}) is three terms: motor thrust applied
along world $+z$ regardless of any launch angle (\src{model/RocketModel.cpp}{70}), gravity
evaluated at the integrator's \emph{trial} position so each stage sees a consistent state
(\src{model/RocketModel.cpp}{77}), and quadratic drag with a single user-set scalar $C_d$ and
reference area (defaults $C_d=1.0$, $A=1.134\times10^{-3}\unit{m^2}$;
\src{model/RocketModel.h}{155}):

\begin{minted}{cpp}
// Clamp altitude to >= 0: trial states crossing z=0 fall outside the atmosphere models' domain.
const double altitude = position[2] > 0.0 ? position[2] : 0.0;
const double rho = atmosphere->getDensity(altitude);
const double speed = velocity.norm();
const Vector3 drag = -0.5 * rho * speed * dragCoefficient * referenceArea * velocity;
\end{minted}

\noindent The $|v|\cdot v$ formulation gives zero drag at $v=0$ with no division, and the altitude
clamp exists because the table-backed atmosphere throws \cppinline|std::out_of_range| below its
lowest bin (\src{model/RocketModel.cpp}{84}, \src{utils/Bin.cpp}{51}). What this force model does
\emph{not} contain --- angle of attack, normal force, a stability moment, the Barrowman machinery
that the model layer already composes but the force path never reads
(\src{model/Propagatable.h}{61}) --- is the subject of \cref{sec:aero}.

The six-degree-of-freedom scaffolding deserves explicit accounting because it is easy to mistake
for capability. \code{StateData} carries orientation and orientation-rate quaternions (initialized
to the degenerate all-zero quaternion, not identity), a DCM, and 3-2-1 Euler angles
(\src{sim/StateData.h}{46}); a repository-wide search finds no code that writes or reads any of
those fields. The orientation integrator is a commented-out member with a design note about the
future \cppinline|DESolver<Quaternion>| ODE (\src{sim/Propagator.h}{103}), and
\code{RocketModel::getTorques} returns \cppinline|Vector3{0.0, 0.0, 0.0}|
(\src{model/RocketModel.cpp}{94}). One piece of the seam is already live: every recorded step
snapshots composite mass, CG, and inertia via \code{writeMassProperties}
(\src{model/RocketModel.cpp}{53}), with the header naming these fields ``the seam the 6-DOF
rotational ODE will read'' (\src{sim/StateData.h}{59}). Attitude propagation is therefore
\fabsent{} at commit \code{20eca72}, but the interfaces it will need are in place and typed.

\subsection{The run loop and typed termination}\label{sec:sim-runloop}

\tikzfig{propagator-loop}{The \code{runUntilTerminate} flight loop
(\srccap{sim/Propagator.cpp}{58}): one integrator step, a finiteness check, the launch-pad hold,
state recording, and the guard cascade that maps every exit onto one of the five typed termination
reasons. The \code{IntegratorError} path is the loop-scope catch that converts an RKF45 step-size
underflow throw into a verdict.}{fig:sim-loop}

\code{Propagator::runUntilTerminate} (\src{sim/Propagator.cpp}{58}) begins every run clean --- a
fresh \code{Nominal} verdict and zeroed statistics (\src{sim/Propagator.cpp}{63}) --- then loops:
read the current state, take one integrator step, and run a guard cascade
(\cref{fig:sim-loop}). The first guard rejects a non-finite state \emph{before} it is recorded or
read by the termination predicate, because a NaN coordinate fails every comparison and the $z<0$
test would never trip (\src{sim/Propagator.cpp}{94}). The second is the launch-pad hold: thrust
curves begin with an implicit $(0,0)$ sample, so over the first step a from-rest rocket sees
essentially zero thrust, accelerates downward, and would cross $z<0$ and terminate after one step.
Until the rocket first achieves $z>0$ under power it is therefore clamped to $z=0$ at rest ---
physically, the pad's normal force cancels the downward force
(\src{sim/Propagator.cpp}{105}). A rocket that can never lift stays put with
$\text{maxAltitude}=0$, which is what arms the \code{NoLiftoff} guard instead of letting the run
end silently.

The five \code{TerminationReason} values (\src{sim/Propagator.h}{35}) partition every exit:

\begin{description}
  \item[\code{Nominal}] the model's own predicate fired. \code{RocketModel} defines end of flight
    as \emph{descending below the launch site}: \cppinline|position[2] < 0.0 && velocity[2] < 0.0|
    (\src{model/RocketModel.cpp}{62}). There are no recovery or staging events; every flight ends
    ballistic at that crossing.
  \item[\code{NoLiftoff}] more than $3\unit{s}$ elapsed without the peak altitude exceeding
    $1\unit{m}$ (thresholds at \src{sim/Propagator.h}{26}; guard at
    \src{sim/Propagator.cpp}{146}) --- an underpowered motor or a hover-balance flight.
  \item[\code{NonFiniteState}] NaN/Inf state caught before recording, so the poisoned sample never
    enters the history (\src{sim/Propagator.cpp}{97}).
  \item[\code{MaxSimTimeExceeded}] the hard backstops: simulated time past \code{maxSimTime}
    (default $7200\unit{s}$, settable; \src{sim/Propagator.h}{115}) \emph{or} $10^{8}$ iterations
    (\src{sim/Propagator.h}{28}).
  \item[\code{IntegratorError}] the whole loop is wrapped in try/catch, so an RKF45 step-size
    underflow throw becomes a verdict instead of escaping (\src{sim/Propagator.cpp}{81},
    \src{sim/Propagator.cpp}{169}).
\end{description}

\begin{hardfail}[Iteration-cap abort is misreported]
The sim-time cap and the $10^{8}$-iteration cap share one branch and one verdict:
\cppinline{if(currentTime > maxSimTime || ++iterations >= maxIterations)} sets
\code{MaxSimTimeExceeded} (\src{sim/Propagator.cpp}{157}), and the enum has no distinct
iteration-cap value (\src{sim/Propagator.h}{40}). A run killed by adaptive step collapse ---
millions of vanishingly small steps at low simulated time --- reports a time-cap exceedance even
when simulated time is nowhere near \code{maxSimTime}. The abort is correct; the diagnosis is
wrong. \fail
\end{hardfail}

Recording is gated by a \code{saveStates} flag: when saving, each accepted step snapshots mass
properties and appends \code{(time, StateData)} to the history stored on the \code{Propagatable}
(\src{sim/Propagator.cpp}{128}, \src{model/Propagatable.h}{47}). Two details matter downstream.
First, the timestamp convention lags one step: \cppinline|appendState(currentTime, nextState)|
records the state reached at $t+h$ against timestamp $t$, because \cppinline|currentTime| advances
by the step actually taken only \emph{after} recording (\src{sim/Propagator.cpp}{133},
\src{sim/Propagator.cpp}{166}). Every consumer of the history --- plots, CSV export, apogee times
--- inherits the offset; a propagator test widens its apogee-time tolerance to $3\,dt$ and names
``the recorder's one-step time-stamp offset'' rather than pinning the convention as intended
(\src{tests/PropagatorTests.cpp}{181}). Second, the \code{saveStates=false} path the loop
carefully supports is dead: \code{retainStates} has no callers anywhere
(\src{sim/Propagator.h}{62}).

\code{TrajectoryStatistics} is updated on every step regardless of \code{saveStates}, so the
flight summary and the no-liftoff guard work even without a retained history
(\src{sim/Propagator.cpp}{127}). The fold uses strict $>$ comparisons --- first attainment wins,
and NaN samples are ignored by construction --- tracking peak altitude and its time, peak speed
($|v|$ norm) and its time, and total flight time as the last sample time
(\src{sim/TrajectoryStatistics.h}{26}).

\subsection{The Integrator facade and the two solver backends}\label{sec:sim-solvers}

Both solvers derive from \cppinline|DESolver<T>| (with \code{T} restricted to \code{Vector3} or
\code{Quaternion} by \cppinline|static_assert|; \src{sim/RK4Solver.h}{33},
\src{sim/RK45Solver.h}{51}) and advance the coupled (state, rate) pair, returning a
\cppinline|StepResult{state, rate, stepSize}| so the caller advances its clock by the step the
solver \emph{actually} took (\src{sim/DESolver.h}{21}). The \code{Quaternion} instantiation is
never used --- another finished seam with zero consumers.

\code{RK4Solver} is textbook fixed-step RK4 with stages at $t$, $t+dt/2$, $t+dt/2$, $t+dt$
(\src{sim/RK4Solver.h}{45}); its timestep starts as NaN, and calling \code{step()} before
\code{setTimeStep()} is a logged error returning a zero result rather than undefined behavior
(\src{sim/RK4Solver.h}{49}). \code{RK45Solver} implements the classical Fehlberg tableau --- six
stages with the Butcher coefficients as \cppinline|constexpr| arrays
(\src{sim/RK45Solver.h}{169}), each stage evaluated at its node time $t + C_n h$
(\src{sim/RK45Solver.h}{92}) so a thrust transient is visible to the error estimate. The
controller is a single absolute norm over state and rate,
$\mathrm{err}=\sqrt{\lVert y_5 - y_4\rVert^2_{\text{state}}+\lVert y_5 - y_4\rVert^2_{\text{rate}}}$
(\src{sim/RK45Solver.h}{120}): a rejected step shrinks $h$ by $(\mathrm{tol}/\mathrm{err})^{1/4}$
floored at $0.1\times$ (\src{sim/RK45Solver.h}{124}); an accepted step grows the next guess by
$(\mathrm{tol}/\mathrm{err})^{1/5}$ capped at $5\times$ (\src{sim/RK45Solver.h}{132}); below
$h_{\min}=10^{-15}$ it throws (\src{sim/RK45Solver.h}{80}); and the fifth-order estimate is
returned as the result (\src{sim/RK45Solver.h}{150}).

\begin{keyinsight}[The vacuum step-runaway clamp]
The next-step guess is clamped to $h_{\max}$, by default $10\times$ the seeded timestep
(\src{sim/RK45Solver.h}{144}, \src{sim/RK45Solver.h}{165}). This is not a tuning nicety but a
correctness requirement peculiar to embedded-pair steppers: constant-acceleration vacuum coasting
is integrated \emph{exactly} by both embedded orders, so $\mathrm{err}\approx 0$ every step, the
$5\times$ growth applies indefinitely, and the integrator leaps past apogee and the ground in a
few giant steps. The failure was real and is now regression-guarded
(\src{sim/tests/RK45SolverTests.cpp}{75}).
\end{keyinsight}

The adaptive solver's tuning surface stops at its own header. \code{setErrorTolerance} and
\code{setMaxStepSize} exist on the concrete class (\src{sim/RK45Solver.h}{70}), but
\code{Integrator} stores its solvers as \cppinline|unique_ptr<DESolver<Vector3>>|
(\src{sim/Integrator.h}{86}) and exposes neither, and neither do \code{Propagator},
\code{QtRocket}, the GUI, or the CLI. Users get the hardcoded $10^{-6}$ absolute tolerance
(\src{sim/RK45Solver.h}{46}) and the $10\times$-seed $h_{\max}$, full stop; the header itself
flags the unscaled absolute norm as ``a possible refinement'' (\src{sim/RK45Solver.h}{36}).

Selection is a string-keyed facade: \code{Integrator} maps \code{Runge-Kutta 4th Order} (the
default) and \code{Runge-Kutta-Fehlberg} to freshly constructed solvers on each
\code{setIntegratorModel} call, and an unknown name is a \emph{logged} no-op that keeps the
current model usable (\src{sim/Integrator.h}{52}). The name travels from the GUI's integrator
combo (populated from \code{getAvailableIntegratorModels()};
\src{gui/SimOptionsTab.cpp}{74}) or the CLI \code{setintegrator} command
(\src{cli/Repl.cpp}{714}) through \code{QtRocket::setIntegratorModel}
(\src{core/QtRocket.h}{36}) to \code{Propagator::setIntegratorModel} (\src{sim/Propagator.h}{90}). The
facade has one sharp edge: \code{setIntegratorFunction} only stores the lambda
(\src{sim/Integrator.h}{72}), so the active solver keeps whatever function it captured at
construction --- the constructor even builds its default RK4 with a null ODE function
(\src{sim/Integrator.h}{28}) --- and a freshly selected solver needs its timestep re-seeded. The
required call order is documented in a test rather than enforced by the API
(\src{sim/tests/IntegratorTests.cpp}{36}); \code{Propagator} gets it right and re-asserts the
configured step at the top of every run (\src{sim/Propagator.cpp}{60}).

\subsection{The Environment registry}\label{sec:sim-env}

\code{Environment} is a string-keyed registry constructed with \code{Constant Gravity} and
\code{Constant Atmosphere} defaults (\src{sim/Environment.h}{31}). All of its models are genuinely
runtime-selectable from both front ends --- GUI combos populated from the registries
(\src{gui/SimOptionsTab.cpp}{72}) and the CLI \code{setatmosphere}/\code{setgravity} commands
(\src{cli/Repl.cpp}{658}, \src{cli/Repl.cpp}{685}). \Cref{tab:sim-envmodels} summarizes the
inventory.

\begin{table}[htbp]
\centering\small
\begin{tabularx}{\linewidth}{@{}llXl@{}}
\toprule
Family & Registry key & Behavior & Status \\
\midrule
Atmosphere & \code{Constant Atmosphere} & sea-level ISA constants at every altitude
  ($\rho=1.225$, $p=101325$, $T=288.15$; \srccap{sim/ConstantAtmosphere.h}{15}) & \fpresent \\
Atmosphere & \code{US Standard 1976} & seven-layer USSA76 closed forms, 0--80$\unit{km}$
  validated (\srccap{sim/USStandardAtmosphere.cpp}{93}) & \fpresent \\
Atmosphere & \code{Vacuum} & all properties zero; the drag-free baseline
  (\srccap{sim/VacuumAtmosphere.h}{17}) & \fpresent \\
Gravity & \code{Constant Gravity} & $(0,0,-9.80665)$ (\srccap{sim/ConstantGravityModel.h}{18})
  & \fpresent \\
Gravity & \code{Spherical Gravity} & Newtonian inverse-square over the geoid ground radius
  (\srccap{sim/SphericalGravityModel.cpp}{28}) & \fpresent \\
Geoid & (fixed member) & WGS84 mean radius for any lat/lon
  (\srccap{sim/SphericalGeoidModel.cpp}{18}) & \fpartial \\
Wind & (not registered) & zero-vector stub, zero call sites (\srccap{sim/WindModel.cpp}{16})
  & \fabsent \\
\bottomrule
\end{tabularx}
\caption{The \code{Environment} inventory at commit \code{20eca72}. The three atmospheres and two
gravity models are live and selectable from both front ends; the geoid is a single hardwired
member; the wind class is compiled but unreachable.}
\label{tab:sim-envmodels}
\end{table}

The US Standard Atmosphere is the core's showpiece of physical fidelity: seven layers with bases
at 0/11/20/32/47/51/71$\unit{km}$ held in static \code{utils::Bin} tables of base temperature,
density, pressure, and lapse rate (\src{sim/USStandardAtmosphere.cpp}{21}), evaluated in closed
form within each layer --- the exponential barometric formula in isothermal layers and the
power-law form with exponent $gM/(R^{*}L)$ in gradient layers
(\src{sim/USStandardAtmosphere.cpp}{93}) --- rather than by table interpolation. Density,
pressure, and temperature are validated against published NOAA report values to 0.1\% (most rows)
and 0.5\% (a few density rows) up to $80\unit{km}$ (\srcf{sim/tests/USStandardAtmosphereTests.cpp}).
Speed of sound ($\sqrt{\gamma R^{*}T/M}$) and Sutherland viscosity are implemented
(\src{sim/USStandardAtmosphere.cpp}{143}) but sit unconsumed: the force path samples only density,
and no test asserts either property for this model --- the only speed-of-sound assertions anywhere
target the Constant and Vacuum models. The edges are soft in both directions: below $0\unit{m}$
the \code{Bin} lookup throws (guarded by the force-path altitude clamp;
\src{utils/Bin.cpp}{51}), while above $71\unit{km}$ the last layer's formula extrapolates
indefinitely because falling off the table's end returns the last bin
(\src{utils/Bin.cpp}{57}) --- the real USSA76 regime change above $80\unit{km}$ is not modeled,
and the header still carries an unresolved verification TODO
(\src{sim/USStandardAtmosphere.h}{18}).

Spherical gravity is a true Newtonian inverse square: the launch frame maps to geocentric
coordinates as $(x, y, z + \text{groundLevel})$ and $a = -GM\,\vec{r}/|\vec{r}|^{3}$ with
$GM=3.986004418\times10^{14}$ (\src{sim/SphericalGravityModel.cpp}{28},
\src{utils/math/Constants.h}{26}). Its geoid, however, is a placeholder twice over: the single
\code{SphericalGeoidModel} accepts latitude and longitude and ignores both, returning the WGS84
mean radius $6371008.8\unit{m}$ unconditionally (\src{sim/SphericalGeoidModel.cpp}{18},
\src{utils/math/Constants.h}{24}), and the gravity model queries it once at construction with
literal $(0,0)$ --- the comment reads ``(0, 0) is a placeholder until a real launch site is
plumbed in'' (\src{sim/SphericalGravityModel.cpp}{19}). There is no launch-site concept anywhere
in \code{Environment} or the controller, so altitude above the pad is implicitly altitude above
sea level in every atmosphere lookup.

One registry behavior is inconsistent with the integrator facade and worth stating plainly:
\code{setGravityModel} and \code{setAtmosphereModel} have no else branch, so an unrecognized name
is a \emph{silent} no-op (\src{sim/Environment.h}{60}) --- tests pin the no-op
(\src{sim/tests/EnvironmentTests.cpp}{94}) --- whereas the same mistake on
\code{setIntegratorModel} is logged as an error (\src{sim/Integrator.h}{64}). Today both front ends
populate their choices from the registries themselves, so a typo cannot arise; a future config
file or save format would make it invisible.

Wind is the purest specimen of the seam-without-consumer motif. The class exists and compiles into
the sim library (\src{core/sim/CMakeLists.txt}{24}); its entire behavior is

\begin{minted}{cpp}
Vector3 WindModel::getWindSpeed(double /* x */, double /* y */ , double /* z */)
{
   return Vector3(0.0, 0.0, 0.0);
}
\end{minted}

\noindent (\src{sim/WindModel.cpp}{16}), and a repository-wide search finds zero call sites: it is
never instantiated, has no slot in the \code{Environment} registry
(\src{sim/Environment.h}{103}), and never enters \code{getForces} --- drag uses ground-relative
velocity only (\src{model/RocketModel.cpp}{88}). Wind cannot affect any simulation at this
commit; what that costs in flight fidelity is discussed in \cref{sec:aero}.

\subsection{The QtRocket controller}\label{sec:sim-controller}

\code{QtRocket} is the Qt-free controller singleton whose ownership, API surface, and racy
double-checked initialization are examined in \cref{sec:overview}; what matters here is how it
runs a flight. \code{launchRocket()} (\src{core/QtRocket.cpp}{54}) runs the four launch steps of
\cref{sec:overview} synchronously, end to end, on the calling thread
(\code{RocketModel::launch} copies the initial state to current and ignites at $t=0$;
\src{model/RocketModel.cpp}{105}). There is no worker thread, no progress reporting, and no cancellation:
the simulation runs to completion on whichever thread calls it --- the Qt event thread when the
GUI's launch button fires (\src{gui/CannonballTab.cpp}{128}), so the UI freezes for the run's
duration, and the REPL loop for the CLI (\src{cli/Repl.cpp}{755}). Nothing reads the state history
concurrently today and no lock protects it --- an invariant that must be revisited the moment the
recovery work of \cref{sec:roadmap} forces simulation onto a worker thread.

Results flow back by const reference: \code{getStates()} exposes the \code{(time, StateData)}
history stored on the \code{Propagatable} (\src{model/Propagatable.h}{47}), alongside
\code{getTrajectoryStatistics()} and \code{getTerminationReason()} (\src{core/QtRocket.h}{48}). Two
loose ends round out the picture. \code{addRocket} --- the API that would swap in a fresh
model/propagator pair --- has no call sites (\src{core/QtRocket.h}{41}), joining
\code{retainStates} in the dead-API column, and \code{setInitialState}'s doc comment still
describes a stale twelve-element pitch/yaw/roll state vector that the 3-DOF core never consumes
(\src{core/QtRocket.h}{66}). Both are consolidated with the other incompleteness findings in
\cref{sec:roadmap}.
